\documentclass{article}
\usepackage[T1]{fontenc}
\usepackage{lmodern}
\usepackage{graphicx}
\usepackage{amsmath,amssymb}
\usepackage[a4paper,margin=2cm]{geometry}
\usepackage[absolute,overlay]{textpos}
\setlength{\TPHorizModule}{1mm}
\setlength{\TPVertModule}{1mm}
\usepackage[indonesian]{babel}
\usepackage{hyperref}
\setlength{\emergencystretch}{2em}
% unit: Halaman 2

\begin{document}

% === Halaman 2 (llm) ===

\section*{Pendahuluan}

\begin{itemize}
    \item Keseluruhan \textit{point} dari kriptografi adalah menjaga kerahasiaan pesan atau kunci dari pihak ketiga. Pihak ketiga: musuh, lawan, penyadap, kriptanalisis (\textit{cryptanalyst}).
    \item Pihak lawan berusaha memecahkan cipherteks dengan melakukan serangan terhadap sistem kriptografi.
    \item Tujuan serangan adalah untuk mengungkap plainteks dari cipherteks atau mendapatkan kunci dekripsi, sehingga kunci dekripsi dapat digunakan untuk mendekripsi potongan cipherteks lainnya.
\end{itemize}

\end{document}
